%! AP Statistics — Unit 2, Lesson 2.9 — Group Activity
\documentclass[10pt]{article}
\usepackage{apstats-article}
\usepackage{apstats-boxes}

\begin{document}

\pageheader{Unit 2, Lesson 2.9}{Group Activity}
\namepartnerperiod

\vspace{0.18in}

% ── SCENARIO ─────────────────────────────────────────────────────────────────
\begin{scenariobox}[Context]{sky}
\small
Each tier works with a real-world dataset that has a nonlinear relationship. Your group will
diagnose the nonlinearity, apply a log transformation, fit a regression to the transformed
data, and back-transform to make a prediction. Show all work and write answers in context.
\end{scenariobox}

\vspace{0.12in}

% ── TIER R ────────────────────────────────────────────────────────────────────
\begin{tcolorbox}[colback=white, colframe=black!40, boxrule=0.8pt, arc=2mm,
    left=4mm, right=4mm, top=3mm, bottom=3mm,
    title={\textbf{Tier R --- Remediate}}, fonttitle=\bfseries, halign title=center]
\small

\textbf{Context: Smartphone Resale Value}

A researcher recorded the age of used smartphones (years, $x$) and their resale value
(\$, $y$). The scatterplot curves downward (exponential decay). A $\log_{10}$ transformation
of $y$ produces a linear scatterplot. Technology output for the regression of $\log y$ on $x$:
\[
  \widehat{\log(\text{value})} = 2.80 - 0.18x \qquad r^2 = 0.94
\]

\vspace{8pt}

\begin{enumerate}[leftmargin=1.6em, itemsep=0.18in]

  \item Which variable was transformed? \blank{3cm} \quad What type of model is suggested?
        \blank{4cm}

  \item Write the back-transformed equation for predicting the original resale value.
        Show work.\\[8pt]
        \writeline\\[3pt]\writeline

  \item Predict the resale value of a smartphone that is \textbf{3 years old}. Show all work.\\[8pt]
        \writeline\\[3pt]\writeline\\[3pt]\writeline

  \item The residual plot for the \textit{original} linear fit showed a curved pattern.
        What does it mean that $r^2 = 0.94$ for the \textit{transformed} model?\\[8pt]
        \writeline\\[3pt]\writeline

\end{enumerate}
\end{tcolorbox}

\vspace{0.12in}

% ── TIER A ────────────────────────────────────────────────────────────────────
\begin{tcolorbox}[colback=white, colframe=black!40, boxrule=0.8pt, arc=2mm,
    left=4mm, right=4mm, top=3mm, bottom=3mm,
    title={\textbf{Tier A --- Approaching Proficiency}}, fonttitle=\bfseries, halign title=center]
\small

\textbf{Context: Kepler's Third Law (Planetary Orbits)}

An astronomer recorded orbital period ($y$, Earth years) and mean orbital radius ($x$,
AU) for eight planets. The scatterplot curves upward. After applying $\log_{10}$ to
\textbf{both} $x$ and $y$, the new scatterplot is nearly perfectly linear. Technology output:
\[
  \widehat{\log(\text{period})} = 0.001 + 1.499\,\log(\text{radius}) \qquad r^2 = 0.9999
\]

\vspace{8pt}

\begin{enumerate}[leftmargin=1.6em, itemsep=0.18in]

  \item Which variable(s) were transformed? \blank{5cm} \quad What type of model?
        \blank{3cm}

  \item Write the regression equation in transformed form (using variable names).\\[6pt]
        \writeline

  \item Write the back-transformed power-model equation. Show work.\\[8pt]
        \writeline\\[3pt]\writeline\\[3pt]\writeline

  \item Predict the orbital period of a planet with a mean radius of \textbf{5.2 AU}.
        Show all work and give the answer in Earth years.\\[8pt]
        \writeline\\[3pt]\writeline\\[3pt]\writeline

  \item The theoretical exponent in Kepler's Third Law is exactly 1.5. Does the data
        support this? Explain using the regression output.\\[8pt]
        \writeline\\[3pt]\writeline

\end{enumerate}
\end{tcolorbox}

\vspace{0.12in}

% ── TIER E ────────────────────────────────────────────────────────────────────
\begin{tcolorbox}[colback=white, colframe=black!40, boxrule=0.8pt, arc=2mm,
    left=4mm, right=4mm, top=3mm, bottom=3mm,
    title={\textbf{Tier E --- Extension}}, fonttitle=\bfseries, halign title=center]
\small

\textbf{Context: Comparing Two Transformations}

A biologist studying fish length ($y$, cm) vs.\ age ($x$, years) has summary statistics for
two candidate models:

\vspace{6pt}
\renewcommand{\arraystretch}{1.4}
\begin{tabularx}{\linewidth}{@{} l X X @{}}
\rowcolor{sky} \textbf{Model} & \textbf{Regression equation} & \textbf{$r^2$} \\
Linear (no transformation) & $\hat{y} = 8.2 + 3.4x$ & 0.78 \\
Transformed ($\log y$ on $x$) & $\widehat{\log y} = 0.92 + 0.11x$ & 0.97 \\
\end{tabularx}

\vspace{10pt}

\begin{enumerate}[leftmargin=1.6em, itemsep=0.18in]

  \item Which model is more appropriate? State \textbf{two} pieces of evidence.\\[8pt]
        \writeline\\[3pt]\writeline\\[3pt]\writeline

  \item Write the back-transformed exponential equation for the better model. Show work.\\[8pt]
        \writeline\\[3pt]\writeline

  \item Predict the length of a fish that is \textbf{7 years old} using the better model.
        Show all steps including the back-transformation.\\[8pt]
        \writeline\\[3pt]\writeline\\[3pt]\writeline

  \item The linear model predicts a length of $8.2 + 3.4(7) = 32.0$ cm.
        The transformed model predicts a different value. Explain why the two models
        give different predictions, and which you trust more.\\[8pt]
        \writeline\\[3pt]\writeline\\[3pt]\writeline

  \item A colleague proposes trying $\log y$ vs.\ $\log x$ instead (power model).
        Explain what additional check you would perform to decide between the exponential
        and power models, and what the residual plot would look like for the better choice.\\[8pt]
        \writeline\\[3pt]\writeline\\[3pt]\writeline

\end{enumerate}
\end{tcolorbox}

\end{document}
